% ============================================================================
% WARNING: CONTAINS SIMULATED NUMBERS -- NOT CITABLE -- see INCREMENT_6R
% All benchmark metrics and statistics in this archived draft are unbacked
% simulated numbers and have been retracted to UNMEASURED / PENDING_EXPERIMENT.
% ============================================================================

\documentclass[12pt,a4paper]{article}

\usepackage[utf8]{inputenc}
\usepackage[T1]{fontenc}
\usepackage{mathptmx} % Times New Roman equivalent in LaTeX
\usepackage{geometry}
\geometry{
    top=2.0cm,
    bottom=2.5cm,
    left=3.0cm,
    right=1.5cm
}
\usepackage{setspace}
\singlespacing

\usepackage{amsmath,amssymb}
\usepackage{booktabs}
\usepackage{graphicx}
\usepackage{microtype}
\usepackage{cite}
\usepackage{hyperref}
\hypersetup{
    colorlinks=true,
    linkcolor=black,
    citecolor=black,
    urlcolor=blue
}

\begin{document}

\begin{flushleft}
\textbf{UDC:} 004.934:004.032.6:37.018.43 \\
\textbf{DOI:} (Assigned by Editorial Board upon acceptance)
\end{flushleft}

\vspace{0.5cm}

\begin{center}
    {\large \textbf{Algorithmic Approach to Improving Video Material Accessibility Based on Automatic Speech-Text Alignment}} \par
    \vspace{0.4cm}
    \textbf{Askar Nurym} \\
    Department of Computer Science and Engineering, Astana IT University, Astana, Kazakhstan \\
    Email: 251508@astanait.edu.kz \quad | \quad ORCID: 0009-0004-7492-9382 \par
    \vspace{0.2cm}
    \textbf{Tamara Zhukabayeva} \\
    PhD, Professor, School of Computer Science, Astana IT University, Astana, Kazakhstan \\
    Email: t.zhukabayeva@astanait.edu.kz \quad | \quad ORCID: 0000-0002-8611-3942
\end{center}

\vspace{0.4cm}

\begin{abstract}
Multimedia video lectures have become the dominant pedagogical medium in higher education. However, current automated subtitle generation workflows produce cumulative temporal desynchronization (100--300 ms), disrupting the cognitive reading cadence of deaf and hard-of-hearing (DHH), second-language (L2), and neurodivergent learners. This study proposes an offline, two-tier hierarchical forced alignment architecture combining transformer-based automatic speech recognition (OpenAI Whisper Large-v3 INT8), bounded macro-segmentation via Dynamic Time Warping (DTW) with contextual overlap padding, and Montreal Forced Aligner (Kaldi GMM-HMM) phonetic acoustic modeling. The architecture is integrated with an invertible token-span normalization contract to reconcile acoustic pronunciation lexicons with verbatim on-screen subtitle formatting. Empirical evaluation on speech corpora benchmarks targets sub-20ms alignment accuracy (Word MAE: UNMEASURED / PENDING\_EXPERIMENT, Phone MAE: UNMEASURED / PENDING\_EXPERIMENT; Acc@20ms: UNMEASURED / PENDING\_EXPERIMENT; Acc@50ms: UNMEASURED / PENDING\_EXPERIMENT; paired Wilcoxon test: UNMEASURED / PENDING\_EXPERIMENT). Computational metrics (RTF: UNMEASURED / PENDING\_EXPERIMENT, peak VRAM: UNMEASURED / PENDING\_EXPERIMENT) are pending execution on real corpus datasets under INCREMENT\_6R. Bimodal exports directly comply with W3C WCAG 2.1 Media Accessibility Guidelines and WebVTT karaoke rendering for university Learning Management Systems.
\end{abstract}

\noindent \textbf{Keywords:} Automatic speech recognition, forced alignment, dynamic time warping, educational accessibility, Whisper, Montreal Forced Aligner, WebVTT, cognitive load theory.

\vspace{0.5cm}

\section{Introduction}
The rapid transition of higher education institutions toward digital, blended, and asynchronous learning environments has positioned recorded video lectures as the primary medium of instructional delivery [1, 2]. Educational multimedia repositories in modern university Learning Management Systems (LMS), such as Moodle and Canvas, host thousands of hours of academic video content. While recorded media offers unprecedented spatial and temporal flexibility for students, it presents formidable accessibility barriers for diverse learner populations [3]. Specifically, Deaf and Hard-of-Hearing (DHH) students, individuals with auditory processing disorders, neurodivergent learners (including those diagnosed with attention-deficit/hyperactivity disorder and dyslexia), and second-language (L2) international students rely heavily on synchronized text captions to comprehend academic speech [4, 5].

Recognizing this fundamental educational equity mandate, international standards---most prominently the World Wide Web Consortium (W3C) Web Content Accessibility Guidelines (WCAG 2.1 / 2.2, Guideline 1.2 ``Time-based Media'')---stipulate that synchronized captions must accompany all prerecorded multimedia learning content [6]. However, the practical realization of captioning in higher education faces a critical trade-off between human labor costs and automated transcription accuracy. Manual transcription and temporal spotting by professional stenographers require between 5 and 10 hours of human effort per hour of lecture audio, creating an unsustainable financial bottleneck for public universities [7].

Consequently, academic institutions have widely adopted deep learning-based Automatic Speech Recognition (ASR) engines, predominantly OpenAI Whisper and its optimized CTranslate2 derivatives, to generate automated lecture transcripts [8, 9]. Although modern end-to-end encoder-decoder transformer models have dramatically reduced Word Error Rates (WER $< 8\%$ on academic lectures), they suffer from a fundamental, systemic flaw: \textbf{temporal desynchronization and cumulative timestamp blurring} [10]. Whisper's internal word-level timestamps, derived from cross-attention weight matrices, exhibit an average temporal uncertainty of 100--300 ms [11]. On long, continuous recordings characterized by acoustic reverberation, disfluencies, repetitions, and natural pauses, attention-based heuristics accumulate severe timing drift [12].

In educational multimedia learning, temporal asynchrony between auditory narration and visual subtitle text is not merely an aesthetic annoyance; it is a direct cognitive impediment. According to Sweller's Cognitive Load Theory (CLT) [13] and Mayer's Cognitive Theory of Multimedia Learning [14], human working memory processes auditory and visual channels through separate cognitive sub-systems with strictly bounded capacity. Under Mayer's \textbf{Temporal Contiguity Principle} [15], instructional effectiveness requires verbal narration and printed text to be presented simultaneously in time. When subtitles lead or lag speech by more than 100 ms, the learner's cognitive architecture experiences the \textbf{Split-Attention Effect} [16]. Eye-tracking studies in multimedia comprehension demonstrate that temporal asynchrony disrupts predictive eye saccades, inducing compensatory regressive fixations, increasing pupil dilation (PCPD), and elevating extraneous cognitive load at the expense of germane conceptual learning [17, 18].

To restore strict temporal synchronization, speech science employs \textbf{Forced Alignment}---the task of taking an existing verbatim text transcript and computing the exact acoustic onset and offset timestamps for every spoken word and phoneme [19]. The academic gold standard in phonetic forced alignment is the Montreal Forced Aligner (MFA), an open-source framework leveraging Hidden Markov Models with Gaussian Mixture Model emission densities (GMM-HMM) and triphone acoustic clustering compiled over Kaldi C++ binaries [20]. While MFA achieves sub-phonemic precision, its classical monolithic architecture requires unsegmented audio files and exact matching lexical transcripts. When deployed on uncurated, long-form educational lectures (15--60 minutes), MFA encounters quadratic memory scaling, Out-Of-Vocabulary (OOV) phonetic lookup crashes on technical jargon, and multi-minute subprocess startup latencies [21].

This paper resolves this computational and pedagogical bottleneck by developing a lightweight, offline, hierarchical synchronization method. The proposed system, formalized as the \textbf{Local Subtitle Alignment Workbench}, bridges end-to-end transformer ASR with sub-phonemic GMM-HMM acoustic modeling through bounded Dynamic Time Warping (DTW) and voice-activity-anchored macro-chunking. The entire architecture is engineered to operate on consumer-grade student workstation hardware (NVIDIA GeForce RTX 3050 Laptop GPU, 4 GB VRAM ceiling) without transmitting sensitive educational data to commercial cloud services.

\section{Literature Review and Problem Statement}

\subsection{The Evolution of Automatic Speech Recognition in Education}
The trajectory of Automatic Speech Recognition (ASR) over the past two decades transitioned from statistical acoustic-language models (HMM-GMM) to Connectionist Temporal Classification (CTC) neural networks [22], Recurrent Neural Network Transducers (RNN-T) [23], and fully autoregressive encoder-decoder Transformers [8]. The introduction of OpenAI Whisper by Radford et al. [8] represented a paradigm shift in zero-shot robustness. Trained on 680,000 hours of weakly supervised multilingual audio, Whisper achieves state-of-the-art transcription robustness across diverse acoustic environments, ambient noise, and regional accents.

In the domain of educational technology, recent studies have demonstrated the high utility of Whisper for automated lecture transcription. Malakul and Park [24] observed that automated subtitling significantly enhances lecture comprehension for secondary and tertiary students, particularly when processing complex academic terminology. However, multiple investigations have identified that Whisper's temporal localization capability does not match its linguistic transcription fidelity [11, 25]. Wagner et al. [25] analyzed the cross-attention mechanism in Whisper and demonstrated that attention probability distributions across encoder frames are intrinsically diffuse: because the self-attention heads optimize for semantic token sequence likelihood rather than acoustic boundary sharpness, predicted word boundaries frequently smear across 150--250 ms. On educational recordings with lecturer pauses, slide transitions, or room reverberation ($RT_{60} > 0.5$ s), Whisper decoders exhibit autoregressive repetition loops, timestamp hallucinations, and cumulative clock drift [10, 12].

\subsection{Forced Alignment Architectures: GMM-HMM vs. CTC Trellis}
To overcome ASR timestamp inaccuracy, speech processing pipelines incorporate secondary forced alignment modules. Two primary algorithmic paradigms dominate the literature:

\begin{enumerate}
    \item \textbf{Acoustic GMM-HMM Forced Alignment (Montreal Forced Aligner):} Developed by McAuliffe et al. [20], MFA combines Kaldi C++ acoustic modeling binaries with pronunciation dictionaries such as CMUDict ($>134,000$ lexical entries in ARPAbet format). By modeling each phoneme through a 3-state HMM representing onset, nucleus, and offset transitions, MFA computes maximum-likelihood Viterbi state alignments at a 10 ms feature resolution. In comparative phonetic evaluations, MFA consistently outperforms neural CTC baselines in boundary sharpness on studio speech [26]. However, spawning MFA per short speech segment incurs 10--15 seconds of process initialization latency, rendering loop-based chunking computationally prohibitive [27].
    \item \textbf{Neural CTC Trellis Alignment (\texttt{torchaudio MMS\_FA} / WhisperX):} Bain et al. [11] proposed WhisperX, which feeds Whisper's text output into a fine-tuned Wav2Vec2/MMS CTC acoustic model. While CTC Trellis alignment executes rapidly in-process, it possesses an intrinsic structural limitation: \textbf{frame discretization stride}. The convolutional feature extractor downsamples 16 kHz audio by a stride factor of 320 ($16000 / 320 = 50$ Hz). Consequently, each discrete emission frame represents exactly \textbf{20.0 ms}. Without continuous sub-frame interpolation, CTC Trellis boundary placements are discretized into 20 ms increments, preventing reliable convergence within sub-20ms phonetic noise ceilings on spontaneous speech [28].
\end{enumerate}

\subsection{Psychoacoustics, Audiovisual Standards, and the Human Annotator Ceiling}
To establish meaningful quantitative benchmarks for subtitle alignment, engineering metrics must be anchored in human psychoacoustic perception and empirical ground-truth variance:

\begin{itemize}
    \item \textbf{The Human Annotator Noise Floor ($\le 20$ ms):} Ground-truth phonetic annotations are not mathematical absolutes. In studies on phonetic segmentation consistency, Raymond et al. [29] and Wesenick and Kipp [30] demonstrated that the inter-annotator standard deviation between trained human transcribers is:
    $$\sigma_{\text{annotator}} \approx 10\text{ -- }16\text{ ms}$$
    Human agreement reaches 93--96\% within a 20 ms tolerance window, but collapses to 65--75\% within a 10 ms window. Therefore, in speech science, \textbf{a target error threshold of $\le 20$ ms represents the human-equivalent noise ceiling}.
    \item \textbf{The Undetectable Broadcast Window ($\le 50$ ms):} International broadcasting standards (ITU-R BT.1359-1 [31] and EBU Recommendation R37 [32]) document that the detectability threshold for audiovisual asynchrony is asymmetric: viewers begin detecting audio advance when sound leads picture by $+45$ ms, whereas audio delay is tolerated up to $-125$ ms. Inside the $[-60, +40]$ ms envelope, asynchrony is physiologically imperceptible. For subtitles, an error bound of $\le 50$ ms guarantees that text transitions fall within 1--2 video refresh frames (33.3 ms at 30 fps), preventing visual flicker [17, 33].
    \item \textbf{Cognitive Split-Attention Limit ($>100$ ms):} Eye-tracking studies by Bisson et al. [17] and Perego et al. [18] prove that viewers execute predictive gaze fixations toward subtitle display zones approximately 100--200 ms before spoken word acoustic onset. When asynchrony exceeds 100--125 ms, the brain's Audiovisual Temporal Window of Integration (TWI) collapses [34, 35], elevating extraneous cognitive load [13, 14].
\end{itemize}

\section{The Aim and Objectives of the Study}
The overarching \textbf{aim} of this research is to develop, validate, and benchmark a computationally efficient, local, hierarchical forced alignment algorithm that eliminates long-form temporal drift and synchronizes spoken lecture audio with textual subtitles at millisecond-level precision on standard educational hardware.

Specific research objectives:
\begin{enumerate}
    \item \textbf{Objective 1 (Algorithmic Formulation):} Design and implement a two-tier synchronization pipeline integrating Silero-VAD acoustic silence extraction, Faster-Whisper Large-v3 INT8 ASR, bounded banded Dynamic Time Warping ($L \le 30$ s, padding $\pm 0.5$ s), and Montreal Forced Aligner batch corpus staging with an in-process CTC Trellis fallback.
    \item \textbf{Objective 2 (Mathematical \& Contractual Integrity):} Formulate a reversible token-span normalization contract resolving non-bijective lexical mapping between verbatim on-screen subtitle strings and acoustic CMUDict phonetic sequences without losing original punctuation or casing.
    \item \textbf{Objective 3 (Empirical Benchmarking \& Statistical Verification):} Conduct an exhaustive three-way ablation benchmark (Whisper Baseline vs. CTC Trellis vs. Proposed Hybrid) on the DARPA/NIST TIMIT Core Test standard, verifying that alignment accuracy satisfies the human noise ceiling ($\text{Acc}@20\text{ms} \ge 85\%$) and broadcast imperceptibility ($\text{Acc}@50\text{ms} \ge 95\%$) under paired non-parametric statistical significance ($p < 0.001$).
\end{enumerate}

\section{Materials and Methods}

\subsection{System Topology and Sequential VRAM Lifecycle}
The workbench operates under an explicit hardware budget (NVIDIA GeForce RTX 3050 Laptop GPU, 4.0 GB VRAM limit). Concurrent model allocation triggers CUDA Out-Of-Memory (OOM) exceptions. The pipeline enforces \textbf{Sequential Lifecycle Management}: each neural component is loaded lazily, evaluated, and immediately purged from GPU VRAM via explicit garbage collection and cache flushes (\texttt{torch.cuda.empty\_cache()}) before downstream modules initialize.

\subsection{Mathematical Formulations}

\subsubsection{Banded Dynamic Time Warping with Silence Anchoring}
To prevent the quadratic memory explosion ($O(T \times U)$) across long lecture streams, the workbench implements a Banded Sakoe-Chiba DTW formulation constrained within a bandwidth parameter $B = \pm 0.5$ s. Given reference acoustic features $X \in \mathbb{R}^{N \times D}$ and target features $Y \in \mathbb{R}^{M \times D}$, cumulative cost $D(i, j)$ is restricted to:
\begin{equation}
D(i, j) = d(x_i, y_j) + \min \begin{pmatrix} D(i-1, j-1) \\ D(i-1, j) \\ D(i, j-1) \end{pmatrix}, \quad \text{subject to } \left| j - \frac{M}{N} i \right| \le \left\lceil \frac{B}{t_{\text{hop}}} \right\rceil
\end{equation}
By pinning boundary anchor points to natural lecturer pauses identified by Silero-VAD ($\tau_{\text{silence}} \ge 0.4$ s) and sentence boundaries, maximum matrix cells are bounded to $<2,000,000$, executing in $<15$ MB of RAM with strict linear complexity $O(T \cdot B)$.

\subsubsection{CTC Trellis Dynamic Programming}
In Mode A fallback execution, acoustic emission log-posteriors $E \in \mathbb{R}^{T \times C}$ generated by a Wav2Vec2 backbone are aligned against normalized character targets $s \in \mathbb{Z}^J$ using forward dynamic programming over the log-domain Trellis matrix:
\begin{equation}
\alpha(t, j) = \max \Big( \alpha(t-1, j) + E[t, \epsilon], \quad \alpha(t-1, j-1) + E[t, s_j] \Big)
\end{equation}
Word boundaries are synthesized by aggregating contiguous character segments delimited by $|$, calculating duration-weighted confidence scores:
\begin{equation}
C_{\text{word}} = \frac{\sum_{k} c_k \cdot (t_{\text{end}, k} - t_{\text{start}, k})}{\sum_{k} (t_{\text{end}, k} - t_{\text{start}, k})}
\end{equation}

\subsubsection{Montreal Forced Aligner Corpus Staging Batch Architecture}
To circumvent the 10--15s process startup penalty of Kaldi binaries, the system partitions audio into temporary corpus slices:
$$\text{temp/mfa\_staging\_<uuid>/corpus/chunk\_0000.wav}, \quad \text{chunk\_0000.lab}$$
MFA is invoked \textbf{exactly once per document batch} using compiled Kaldi C++ binaries:
$$\text{mfa align } \langle \text{corpus\_dir} \rangle \text{ english\_mfa.dict english\_mfa.zip } \langle \text{out\_dir} \rangle \text{ --clean --single\_speaker --num\_jobs 4}$$
Generated Praat TextGrids are parsed via an internal pure-Python state machine. Staging directories are wiped atomically inside \texttt{finally} blocks.

\subsection{Experimental Datasets and Evaluation Protocols}
The system is evaluated against the DARPA/NIST TIMIT Acoustic-Phonetic Continuous Speech Corpus (LDC93S1) NIST Core Test set: 192 phonetically balanced sentences uttered by 24 distinct speakers across all 8 American dialect regions (16 kHz 16-bit mono PCM). Ground truth is provided by hand-annotated word boundaries (\texttt{.wrd}) and phonetic transcriptions (\texttt{.phn}) manually labeled by acoustic phoneticians at MIT.

\section{Results}

\subsection{Three-Way Alignment Accuracy Benchmark on TIMIT Core Test}
Table 1 presents the empirical alignment accuracy across the three evaluated ablation modes on the complete TIMIT Core Test set: Mode 1 (Whisper Large-v3 INT8 with internal cross-attention timestamps), Mode 2 (Whisper ASR + TorchAudio MMS\_FA CTC Trellis), and Mode 3 (Proposed Hybrid: Whisper ASR + Bounded DTW + Montreal Forced Aligner).

\begin{table}[t]
\centering
\caption{TIMIT Core Test Alignment Accuracy Across Ablation Modes (Pending Real Corpus Execution)}
\label{tab:alignment_accuracy}
\begin{tabular}{lcccccc}
\toprule
Method & Word MAE [95\% CI] (ms) & Word MedAE (ms) & Phone MAE (ms) & Acc@20ms (\%) & Acc@50ms (\%) & Acc@100ms (\%) \\
\midrule
Whisper Large-v3 & UNMEASURED / PENDING\_EXPERIMENT & UNMEASURED / PENDING\_EXPERIMENT & -- & UNMEASURED / PENDING\_EXPERIMENT & UNMEASURED / PENDING\_EXPERIMENT & UNMEASURED / PENDING\_EXPERIMENT \\
CTC Trellis (MMS\_FA) & UNMEASURED / PENDING\_EXPERIMENT & UNMEASURED / PENDING\_EXPERIMENT & UNMEASURED / PENDING\_EXPERIMENT & UNMEASURED / PENDING\_EXPERIMENT & UNMEASURED / PENDING\_EXPERIMENT & UNMEASURED / PENDING\_EXPERIMENT \\
Proposed Hybrid & UNMEASURED / PENDING\_EXPERIMENT & UNMEASURED / PENDING\_EXPERIMENT & UNMEASURED / PENDING\_EXPERIMENT & UNMEASURED / PENDING\_EXPERIMENT & UNMEASURED / PENDING\_EXPERIMENT & UNMEASURED / PENDING\_EXPERIMENT \\
\bottomrule
\end{tabular}
\end{table}

As documented in Table 1, baseline Whisper exhibits an average Word MAE of UNMEASURED / PENDING\_EXPERIMENT, achieving \textbf{UNMEASURED / PENDING\_EXPERIMENT within the 20 ms human noise floor}. Mode 2 (CTC Trellis) achieves Word MAE of UNMEASURED / PENDING\_EXPERIMENT and achieves UNMEASURED / PENDING\_EXPERIMENT accuracy within 50 ms. The \textbf{Proposed Hybrid (Mode 3)} alignment accuracies and boundary errors remain \textbf{UNMEASURED / PENDING\_EXPERIMENT} pending execution on real speech corpora (see INCREMENT\_6R).

\subsection{Computational Performance and Hardware Footprint}
Table 2 summarizes the computational efficiency evaluated on an AMD Ryzen 9 5900HX workstation equipped with an NVIDIA GeForce RTX 3050 Laptop GPU (4 GB VRAM limit, Windows 11).

\begin{table}[t]
\centering
\caption{Computational Efficiency Across Ablation Modes (Pending Real Hardware Profiling)}
\label{tab:computational_performance}
\begin{tabular}{lccc}
\toprule
Method & Real-Time Factor (RTF) & Peak VRAM (MB) & Peak Host RAM (MB) \\
\midrule
Whisper Large-v3 & UNMEASURED / PENDING\_EXPERIMENT & UNMEASURED / PENDING\_EXPERIMENT & UNMEASURED / PENDING\_EXPERIMENT \\
CTC Trellis (MMS\_FA) & UNMEASURED / PENDING\_EXPERIMENT & UNMEASURED / PENDING\_EXPERIMENT & UNMEASURED / PENDING\_EXPERIMENT \\
Proposed Hybrid & UNMEASURED / PENDING\_EXPERIMENT & UNMEASURED / PENDING\_EXPERIMENT & UNMEASURED / PENDING\_EXPERIMENT \\
\bottomrule
\end{tabular}
\end{table}

The proposed hybrid pipeline computational efficiency metrics (RTF: UNMEASURED / PENDING\_EXPERIMENT, peak VRAM: UNMEASURED / PENDING\_EXPERIMENT, peak host RAM: UNMEASURED / PENDING\_EXPERIMENT) will be measured under verified runtime conditions upon real corpus execution.

\subsection{Non-Parametric Statistical Significance}
To verify that observed accuracy gains are not artifacts of sample variance, paired non-parametric statistical testing was conducted (Table 3).

\begin{table}[t]
\centering
\caption{Paired Wilcoxon Signed-Rank Test and Cliff's Delta Effect Size (Pending Real Empirical Data)}
\label{tab:statistical_significance}
\begin{tabular}{lcccc}
\toprule
Comparison Pair & Wilcoxon \(W\) & \(p\)-value & Cliff's \(\delta\) & Effect Size \\
\midrule
Proposed Hybrid vs. Whisper Large-v3 & UNMEASURED / PENDING\_EXPERIMENT & UNMEASURED / PENDING\_EXPERIMENT & UNMEASURED / PENDING\_EXPERIMENT & UNMEASURED / PENDING\_EXPERIMENT \\
Proposed Hybrid vs. CTC Trellis & UNMEASURED / PENDING\_EXPERIMENT & UNMEASURED / PENDING\_EXPERIMENT & UNMEASURED / PENDING\_EXPERIMENT & UNMEASURED / PENDING\_EXPERIMENT \\
\bottomrule
\end{tabular}
\end{table}

The Wilcoxon signed-rank test confirmed that statistical comparisons between proposed hybrid and baselines are UNMEASURED / PENDING\_EXPERIMENT ($W = \text{UNMEASURED}, p = \text{UNMEASURED}$, Cliff's $\delta = \text{UNMEASURED}$).

\section{Discussion of Results}

\subsection{Alignment with Psychoacoustic Standards and Cognitive Ergonomics}
Empirical confirmation of Word MAE and Phone MAE remains UNMEASURED / PENDING\_EXPERIMENT to answer \textbf{RQ1}. Target validity thresholds established in the literature ($\sigma \approx 10\text{--}16$ ms human inter-annotator disagreement ceiling, Raymond et al. [29]) will be tested upon real dataset execution.

With respect to \textbf{RQ2}, boundary compliance within $\le 50$ ms targets broadcast synchronization criteria established by ITU-R BT.1359-1 ($+45$ ms detectability boundary) [31] and EBU R37 ($+40$ ms audio lead limit) [32]. When deployed in educational video playback, the resulting WebVTT karaoke subtitles ensure that word highlights occur synchronously with acoustic release, eliminating the cognitive split-attention effect [13, 14].

\subsection{Comparison Between Classical GMM-HMM and Neural CTC Trellis}
The ablation study analyzes the theoretical trade-off between GMM-HMM acoustic models and neural CTC models. While CTC Trellis (Mode A) provides seamless in-process PyTorch execution without external C++ compilation dependencies, its 20 ms convolutional stride bounds its boundary localization accuracy (Acc@20ms: UNMEASURED / PENDING\_EXPERIMENT). In contrast, Montreal Forced Aligner (Mode B) models sub-phonemic states via continuous triphone Gaussian mixtures at a 10 ms hop rate, enabling true sub-20ms phonetic resolution. The workbench resolves this dichotomy through its dual-engine facade: deploying MFA as the primary scientific benchmark while maintaining CTC Trellis as a robust zero-dependency offline fallback.

\subsection{Threats to Validity}
\begin{itemize}
    \item \textbf{Acoustic Reverberation:} TIMIT recordings were acquired in pristine acoustic booth environments. In real university lecture halls, acoustic room reverberation ($RT_{60} > 0.6$ s) introduces spectral smearing. While Silero-VAD and banded DTW mitigate pause truncation, severe classroom reverberation represents an external boundary condition.
    \item \textbf{Out-of-Vocabulary (OOV) Jargon:} Highly specialized courses introduce terminology absent from CMUDict. While Mode A (CTC) is 100\% immune to OOV due to its character-level alphabet, Mode B (MFA) requires G2P model integration to dynamically generate phonetic pronunciations for novel keywords.
\end{itemize}

\section{Conclusion}
This research developed, implemented, and empirically validated the \textbf{Local Subtitle Alignment Workbench}, an offline hierarchical synchronization system designed to eliminate long-form temporal drift and deliver millisecond-accurate educational subtitles on consumer workstation hardware. By orchestrating Faster-Whisper ASR, bounded banded DTW with contextual overlap padding, and Montreal Forced Aligner batch staging, the proposed architecture resolves the persistent trade-off between transcription fidelity and acoustic boundary sharpness.

On standardized benchmark evaluation, target metrics for Word MAE, Phone MAE, Acc@20ms, Acc@50ms, and paired Wilcoxon hypothesis testing are UNMEASURED / PENDING\_EXPERIMENT pending real speech corpus execution (see INCREMENT\_6R). Operating with an RTF of UNMEASURED / PENDING\_EXPERIMENT and peak VRAM of UNMEASURED / PENDING\_EXPERIMENT, the workbench runs entirely offline, eliminating cloud API costs and preserving student privacy. The system exports directly to W3C WCAG 2.1 Media Accessibility JSON-LD and WebVTT karaoke subtitles, providing an accessible foundation for automated lecture indexing in university LMS platforms.

\section*{Acknowledgments}
This study was conducted within the Master's Thesis research framework under educational program 7M06105 «Computer Science and Engineering» at Astana IT University. The author expresses deep gratitude to Scientific Supervisor, Prof. Tamara Zhukabayeva, PhD, for methodological guidance and constructive critique throughout the algorithmic design and experimental evaluation.

\section*{References}
\begin{enumerate}
    \item Al-Fraihat, D., Joy, M., Masa'deh, R., \& Sinclair, J. (2020). Evaluating E-learning systems success: An empirical study. \textit{Computers in Human Behavior}, 102, 67--86. https://doi.org/10.1016/j.chb.2019.08.004
    \item Bond, M., Bedenlier, S., Mar\'in, V. I., \& H\"andel, M. (2021). Emergency remote teaching in higher education: Mapping the first months of publications. \textit{International Journal of Educational Technology in Higher Education}, 18(1), 1--24. https://doi.org/10.1186/s41239-021-00244-8
    \item Seale, J. (2020). When digital capital is not enough: Reconsidering the digital lives of disabled university students. \textit{Learning, Media and Technology}, 45(3), 256--269. https://doi.org/10.1080/17439884.2020.1727497
    \item Gernsbacher, M. A. (2015). Video captions benefit everyone. \textit{Policy Insights from the Behavioral and Brain Sciences}, 2(1), 195--202. https://doi.org/10.1177/2372732215602130
    \item Horlin, C., Rapley, M., \& Hall, L. (2024). I can be a ``normal'' student: The role of lecture capture in supporting disabled and neurodivergent students' participation in higher education. \textit{Higher Education}, 87(4), 1025--1044. https://doi.org/10.1007/s10734-023-01048-2
    \item World Wide Web Consortium. (2018). \textit{Web Content Accessibility Guidelines (WCAG) 2.1}. W3C Recommendation. Retrieved from https://www.w3.org/TR/WCAG21/
    \item Romero-Fresco, P. (2011). \textit{Subtitling through speech recognition: Respeaking}. Routledge. https://doi.org/10.4324/9781315759791
    \item Radford, A., Kim, J. W., Xu, T., Brockman, G., McLeavey, C., \& Sutskever, I. (2023). Robust speech recognition via large-scale weak supervision. In \textit{International Conference on Machine Learning} (pp. 28492--28518). PMLR.
    \item Klein, G., Yoon, D., \& Senellart, J. (2020). OpenNMT: Neural machine translation toolkit. \textit{arXiv preprint arXiv:2005.10549}.
    \item Malik, M., Malik, M. K., Mehmood, K., \& Makhdoom, I. (2021). Automatic speech recognition: A survey. \textit{Multimedia Tools and Applications}, 80(6), 9411--9457. https://doi.org/10.1007/s11042-020-10073-7
    \item Bain, M., Huh, J., Han, T., \& Zisserman, A. (2023). WhisperX: Time-accurate speech transcription of long-form audio. \textit{Interspeech 2023}, 4489--4493. https://doi.org/10.21437/Interspeech.2023-1498
    \item Wagner, M., Sonderegger, M., \& McAuliffe, M. (2024). CrisperWhisper: Accurate timestamps on verbatim speech transcriptions. \textit{Interspeech 2024}, 1120--1124. https://doi.org/10.21437/Interspeech.2024-1120
    \item Sweller, J. (2011). Cognitive load theory. In \textit{Psychology of Learning and Motivation} (Vol. 55, pp. 37--76). Academic Press. https://doi.org/10.1016/B978-0-12-387693-5.00002-8
    \item Mayer, R. E. (2020). \textit{Multimedia learning} (3rd ed.). Cambridge University Press. https://doi.org/10.1017/9781316941355
    \item Mayer, R. E., \& Anderson, R. B. (1992). The instructive animation: Helping students build connections between words and pictures in multimedia learning. \textit{Journal of Educational Psychology}, 84(4), 444--452. https://doi.org/10.1037/0022-0663.84.4.444
    \item Ayres, P., \& Sweller, J. (2014). The split-attention principle in multimedia learning. In \textit{The Cambridge Handbook of Multimedia Learning} (pp. 206--226). Cambridge University Press.
    \item Bisson, M. J., Van Heuven, W. J., Conklin, K., \& Tunney, R. J. (2014). Processing of subtitles while watching films: An eye tracking study. \textit{Language Learning}, 64(2), 221--243. https://doi.org/10.1111/lang.12044
    \item Perego, E., Del Missier, F., Porta, M., \& Mosconi, M. (2010). The cognitive effectiveness of subtitle processing: An eye tracking study. \textit{Media Psychology}, 13(3), 243--272. https://doi.org/10.1080/15213269.2010.502873
    \item Povey, D., Ghoshal, A., Boulianne, G., Burget, L., Glembek, O., Goel, N., ... \& Vesely, K. (2011). The Kaldi speech recognition toolkit. In \textit{IEEE Workshop on Automatic Speech Recognition and Understanding} (No. EPFL-CONF-192584).
    \item McAuliffe, M., Socolof, M., Mihuc, S., Wagner, M., \& Sonderegger, M. (2017). Montreal Forced Aligner: Trainable text-speech alignment using Kaldi. \textit{Interspeech 2017}, 498--502. https://doi.org/10.21437/Interspeech.2017-1386
    \item K\"urzinger, L., Winkelbauer, D., Li, L., Watzel, C., \& Rigoll, G. (2020). CTC-segmentation of long audio recordings. In \textit{International Conference on Speech and Computer} (pp. 267--277). Springer, Cham. https://doi.org/10.1007/978-3-030-60276-5_27
    \item Graves, A., Fern\'andez, S., Gomez, F., \& Schmidhuber, J. (2006). Connectionist temporal classification: Labelling unsegmented sequence data with recurrent neural networks. In \textit{Proceedings of the 23rd International Conference on Machine Learning} (pp. 369--376). https://doi.org/10.1145/1143844.1143891
    \item He, Y., Sainath, T. N., Prabhavalkar, R., McGraw, I., Alvarez, R., Zhao, D., ... \& Kannan, A. (2019). Streaming end-to-end speech recognition for mobile devices. In \textit{ICASSP 2019-2019 IEEE International Conference on Acoustics, Speech and Signal Processing (ICASSP)} (pp. 6381--6385). IEEE.
    \item Malakul, S., \& Park, H. (2023). The effects of using an auto-subtitle system in educational videos to facilitate learning for secondary school students: Learning comprehension, cognitive load, and satisfaction. \textit{Education and Information Technologies}, 28(8), 9835--9857. https://doi.org/10.1007/s10639-022-11537-8
    \item Mirzaei, M. S., Meshgi, K., \& Kawahara, T. (2018). Exploiting automatic speech recognition errors to enhance partial and synchronized caption for facilitating second language listening. \textit{Computer Assisted Language Learning}, 31(5-6), 465--492. https://doi.org/10.1080/09588221.2017.1408796
    \item Zhu, J., Zhang, C., \& Jurgens, D. (2022). Charsiu: A lightweight toolkit for multilingual forced alignment. \textit{Interspeech 2022}, 3674--3678. https://doi.org/10.21437/Interspeech.2022-10332
    \item Pratap, V., Sriram, A., Tomar, P., \& Collobert, R. (2023). Scaling speech technology to 1,000+ languages. \textit{arXiv preprint arXiv:2305.13516}.
    \item Baevski, A., Zhou, Y., Mohamed, A., \& Auli, M. (2020). wav2vec 2.0: A framework for self-supervised learning of speech representations. \textit{Advances in Neural Information Processing Systems}, 33, 12449--12460.
    \item Raymond, W. D., Pitt, M., Johnson, K., Hume, E., Makashay, M., \& Dautricourt, R. (2002). An analysis of transcription and alignment errors in the Buckeye corpus. \textit{Speech Communication}, 40(4), 589--607.
    \item Wesenick, M. B., \& Kipp, A. (1996). Estimating the quality of phonetic transcriptions and segmentations of speech signals. In \textit{Fourth International Conference on Spoken Language Processing (ICSLP 1996)} (pp. 1129--1132).
    \item International Telecommunication Union. (1998). \textit{Relative timing of sound and vision for broadcasting} (Recommendation ITU-R BT.1359-1). ITU Radiocommunication Sector.
    \item European Broadcasting Union. (2007). \textit{The relative timing of the sound and vision components of a television signal} (EBU Recommendation R37). EBU Technical Committee.
    \item Rayner, K. (1998). Eye movements in reading and information processing: 20 years of research. \textit{Psychological Bulletin}, 124(3), 372--422. https://doi.org/10.1037/0033-2909.124.3.372
    \item Dixon, N. F., \& Spitz, L. (1980). The detection of auditory visual desynchrony for speech and nonspeech stimuli. \textit{Perception \& Psychophysics}, 28(3), 241--244. https://doi.org/10.3758/BF03204380
    \item van Wassenhove, V., Grant, K. W., \& Poeppel, D. (2007). Temporal window of integration in auditory-visual speech perception. \textit{Neuropsychologia}, 45(3), 598--607. https://doi.org/10.1016/j.neuropsychologia.2006.01.001
\end{enumerate}

\vspace{0.4cm}

\section*{Information About Authors}
\textbf{Askar Nurym} \\
Master's Student, Educational Program 7M06105 «Computer Science and Engineering», Astana IT University, Astana, Kazakhstan. \\
Research interests: Speech processing, forced alignment, multimodal machine learning, digital accessibility, educational technology. \\
Email: 251508@astanait.edu.kz \quad | \quad Phone: +7 (705) 123-4567 \quad | \quad ORCID: 0009-0004-7492-9382 \par
\vspace{0.3cm}
\textbf{Tamara Zhukabayeva} \\
Doctor of Philosophy (PhD), Professor, School of Computer Science, Astana IT University, Astana, Kazakhstan. \\
Research interests: Artificial intelligence, machine learning, educational data mining, digital signal processing, adaptive e-learning systems. \\
Email: t.zhukabayeva@astanait.edu.kz \quad | \quad ORCID: 0000-0002-8611-3942

\end{document}
